A demonstration version with all the tasks may be found at the \href{https://specl.socsci.uci.edu/acquire/data4/iBat5/demo.html}{following link}\footnote{\url{https://specl.socsci.uci.edu/acquire/data4/iBat5/demo.html}}.

\subsection*{Ethics}

All methods employed in this study adhered to relevant guidelines and regulations, such as the Declaration of Helsinki. The Institutional Review Board (IRB) at the University of California, Irvine (UCI) approved the study protocol and experimental procedures under IRB approval number \#3292. Prior to participating in the study, participants provided written informed consent and were made aware that their participation was entirely voluntary and that they could withdraw at any time. 

\subsection*{Planning}

There are two sample sizes to be set in an experiment.  The first is the number of trials per person per task; the second is the number of people.  These sample sizes play different roles.  The first, the number of trials, is paramount as it determines the systematic attenuation and the expected correctness of any subsequent estimate of correlation.  The second, the number of people, is about precision, and it makes little sense to have high precision in a bad estimate.  

\textcite{Rouder.Mehrvarz.2024} provide a means of planning the number of trials needed in a task.  We used their approach, along with the visual illusion pilots they published, to plan our number of trials.  Given the degree of trial noise, 15 trials per person per task was sufficient for a test-retest reliability of over .9.  The number of participants reflects the precision needed to adjudicate between the two main hypotheses that there are no or small correlations across tasks vs. there is a common, one-factor solution.  We settled on 150 individuals because it resulted in correlation estimates that had a worst-case precision of $\pm .06$.  We do not use significance tests here, but if one did, the following power calculation might be informative.  Suppose we wished to discriminate single task-wise correlations of.25 in value from chance with a power of .8 with a one-tailed test (negative correlations are not interpretable and should be taken as evidence for null in significance testing).  This test is well powered at 98 individuals; hence the chosen value of 150 people is reasonable.


\subsection*{Participants}

One-hundred and fifty individuals, recruited through Prolific, served as participants in the experiment. The only constraint was that participants were located in the United States.  Participants were paid  \$5 USD for the 20 minutes it took to complete the battery.

\subsection*{Demographics}

After the data-cleaning process, 138 participants remained. The average age was 35.34 years (SD = 11.01). There were 66 females, 71 males, and 1 participant who preferred not to say.

\subsection*{Tasks} 

The task battery comprises five visual illusions: the Brentano Illusion \parencite{Brentano.1892}, the Ebbinghaus Illusion \parencite[]{Ebbinghaus.1913}, the Poggendorff Illusion \parencite[]{Day.Dickinson.1976}, the Ponzo Illusion \parencite[]{Ponzo.1910}, and the Zöllner Illusion \parencite[]{Judd.Courten.1905}. These illusions are shown in Figure~\ref{fig:illPlots}. In each task, participants adjusted a critical element of the stimulus. In the Brentano illusion (Panels A), participants adjusted the middle chevron to make the blue and red lines appear to be of equal length, with a tendency to set the red line too long and the blue line too short. In the Ebbinghaus Illusion (Panel C), participants adjusted one of the inner circles to match the other in size. In the Poggendorff Illusion (Panel D), participants aligned one of the two disjointed segments (left and right levers) to form a straight line. In the Ponzo Illusion (Panel E), participants adjusted one of the red lines to be the same height as the other red line. Lastly, in the Zöllner Illusion (Panel F), participants adjusted the horizontal lines to appear parallel. The method of adjustment was employed as it is long-standing \parencite{Fechner.1860} and effective in measuring the strength of illusions \parencite[]{Cretenoud.etal.2021}.

For each task, we constructed two versions. The two versions of the Brentano Illusion are shown in Panels A and B. The motivation for using two versions is to control for idiosyncratic response biases. If a participant ran only the version in Panel A and had a natural tendency to set the left segment too large at the expense of the right segment, this tendency would manifest as a larger illusion effect. To control for such a possibility, no matter how remote, we constructed an alternate version so that such a tendency would not greatly impact the true illusion effect. Alternate versions of each illusion were constructed analogously. For the Ebbinghaus Illusion, in Version 1, participants adjusted the right inner circle (surrounded by smaller circles); in Version 2, they adjusted the left circle (surrounded by larger circles). For the Poggendorff Illusion, in Version 1, participants adjusted the right lever vertically. In Version 2, the illusion was flipped vertically, with the downward-facing right lever being adjusted on the right-hand side. For the Ponzo Illusion, in Version 1, participants adjusted the closer red line, and in Version 2, they adjusted the farther red line. Lastly, Zöllner Version 1 is depicted in Panel F of Figure~\ref{fig:ill-tasks}, and Version 2 was the same but mirrored horizontally. We expected the magnitude of an illusion for an individual to be consistent across versions. This expectation held for some illusions, but surprisingly, not for others. As such, we could not collapse across the versions and they feature prominently in the subsequent analyses.


\subsection*{Design}

The design was a $5 \times 2$ within-subject factorial design. The first factor was the task (the specific illusion), and the second factor was the task version.  We chose to repeat each task-by-version combination 15 times for each individual.  This trial size of 15 was chosen as follows.  We ran a pilot experiment to assess the size of illusions, the variability between people, and the variation across replicate trials.  We then used the method described in \textcite{Rouder.Mehrvarz.2024} to obtain an expected reliability of over .95 per illusion.  The method yielded 30 trials for each illusion which we split across versions.  The response measure was how far the adjustment varied from the ground truth.

\subsection*{Procedure}

Participants registered through Prolific, and were then referred to our webserver for the experiment. After reading the consent letter, they were asked to raise the volume on their speakers so they could hear feedback tones and to size their screen so that they could see all stimulus elements.  Participants ran two cycles, and in each cycle, each of the five illusions was presented in a random order. Versions were randomly chosen in the first cycle, and the other version was presented in the second cycle. For each illusion in the cycle, a set of instructions and two practice trials were provided, and then the participant ran 15 replicate trials. Across all tasks and cycles, there were 150 trials in total. Participants had to adjust elements as discussed previously using the arrow keys on their keyboard, and each press moved the critical element a small amount. When they were satisfied with their adjustment, they pressed the spacebar, which advanced to the next trial. Breaks were given after tasks of 15 trials.

We found in a pilot that participants varied greatly in how much time they took before being satisfied with their adjustment. Some would spend tens of seconds going back and forth until they made the spacebar commitment. To keep the experiment fairly regular across individuals, it was necessary to impose an eight-second time limit. Participants could advance to the next trial by pressing the spacebar or wait the eight seconds for automatic advancement.

For each trial in each illusion, we varied a number of parameters. For convenience, several illusions have a \textit{target} size, such as the size of the to-be-matched line segment (Brentano, Ponzo), circumference (Ebbinghaus), or position (Poggendorff). We jittered these target sizes as well as the starting sizes of the adjusted elements. We also jittered other elements, including the overall position on the screen and, in the Zöllner illusion, the length of the main lines and cross lines. This variation prevented the development of stereotypical responses across repeated trials.